% !TeX root = ../教材题库_全部条目_转写工作稿.tex
% Source: XJ-B1-C05-HW-P196-Y1319
% 第5章 / 5.4 函数 y=A sin(ωx+φ) 的图象与性质 / 教材 p.196、197 / PDF 第203、204页
% OCR未校对：公式与文字均可能有误；未完成数学TeX转写；禁止直接发给学生
% 原题及相邻原解截图（核对用）：../题目原图/XJ-B1-C05-HW-P196-Y1319-01.webp, ../题目原图/XJ-B1-C05-HW-P196-Y1319-02.webp
% 本文件只在转写完成后作为题干使用；OCR 条目尚不是可靠题干。
\begin{bankitem}{习题5.4 第2题}{第5章 / 5.4 函数 y=A sin(\ensuremath{\omega}x+\ensuremath{\varphi}) 的图象与性质 / 教材 p.196、197 / PDF 第203、204页}{XJ-B1-C05-HW-P196-Y1319}{OCR未校对：公式与文字均可能有误；未完成数学TeX转写；禁止直接发给学生}
\begin{ocrdraft}
[教材p.196 / PDF第203页]
\par
2．选择题：
\par
《）的图象，丿
\par
《）的图象
\par
（1）为了得到函数3，一2sinx\ocrglyph{00AF}
\par
只需把函数3，=2sin工十
\par
上所有的点（
\par
(A)向左平移匹个单位长度
\par
(B)
\par
向右平移晋个单位长度
\par
（0向左平移个单位长度
\par
向右平移扌个单位长度
\par
[教材p.197 / PDF第204页]
\par
（2）要得到函数3，一cos（2工十1）的图象，只需把函数3，一cos2x的图象上所有的
\par
(A)向左平移1个单位长度
\par
(B)向右平移1个单位长度
\par
(C)向左平移一个单位长度
\par
(D)向右平移个单位长度
\end{ocrdraft}
\end{bankitem}
